\begin{figure}[ht]
\centering
\begin{tikzpicture}[x=1cm, y=1cm, font=\small]

% Illustrate L'Hopital's rule as a ratio of tangent lines near a
% common zero. Both f(x)=sin x and g(x)=x pass through the origin;
% near it each looks like its tangent, and the ratio f/g tends to
% the ratio of slopes f'(0)/g'(0) = 1.
% Plot scale: x by 3.0, y by 3.0, domain x in [-1.05, 1.05].

% Axes
\draw[->, thick] (-3.6, 0) -- (3.6, 0) node[right] {$x$};
\draw[->, thick] (0, -1.1) -- (0, 3.2) node[above] {$y$};

% f(x) = sin x : plotx = 3*x, ploty = 3*sin(x rad)
\draw[very thick, engblue] plot[domain=-1.05:1.05, samples=90]
  ({3*\x}, {3*sin(deg(\x))});
\node[engblue, anchor=south east, font=\scriptsize] at (3.0, {3*sin(deg(1.0))})
  {$f(x) = \sin x$};

% g(x) = x : plotx = 3*x, ploty = 3*x
\draw[very thick, engred] plot[domain=-1.05:1.05, samples=2]
  ({3*\x}, {3*\x});
\node[engred, anchor=south west, font=\scriptsize] at ({3*0.72}, {3*0.9})
  {$g(x) = x$};

% common tangent-line envelope near origin (both slope ~1, coincident)
\fill[engblack] (0, 0) circle (2pt);
\node[engblack, anchor=north east, font=\scriptsize] at (-0.1, -0.08)
  {$f(0) = g(0) = 0$};

% zoom callout: at x = 0.9, show the two ordinates and their ratio
\draw[thin, enggray, dotted] ({3*0.9}, 0) -- ({3*0.9}, {3*0.9});
\fill[engblue] ({3*0.9}, {3*sin(deg(0.9))}) circle (1.4pt);
\fill[engred] ({3*0.9}, {3*0.9}) circle (1.4pt);
\node[anchor=north, font=\scriptsize] at ({3*0.9}, -0.05) {$x$};
\draw[<->, enggray, very thin] ({3*0.9 + 0.18}, {3*sin(deg(0.9))})
  -- ({3*0.9 + 0.18}, {3*0.9});
\node[enggray, anchor=west, font=\tiny] at ({3*0.9 + 0.25}, {3*0.85})
  {gap shrinks as $x \to 0$};

% annotation of the limiting ratio
\node[engamber, anchor=west, font=\scriptsize] at (-3.5, 2.7)
  {$\displaystyle\lim_{x\to 0}\frac{f(x)}{g(x)}
    = \frac{f'(0)}{g'(0)} = 1$};

\end{tikzpicture}
\caption[L'H\^opital's rule as a ratio of linearisations]{The
indeterminate ratio $\sin x / x$ at $x = 0$, where both numerator and
denominator vanish. Near the origin each function is its own tangent
line to first order, $\sin x \approx x$ and $x = x$, so the two graphs
leave the origin along nearly the same direction and the vertical gap
between them shrinks faster than either ordinate. The limit of the
ratio is therefore the ratio of the two slopes at the origin,
$f'(0)/g'(0) = \cos 0 / 1 = 1$, which is what L'H\^opital's rule
computes. The rule is the linearisation of the chapter read as a
quotient: replace each function by its leading linear term and take
the ratio of the coefficients.}
\label{fig:vol02:ch05:lhopital-linearisation}
\end{figure}
